\section{Modèle et notations}\label{sec-model}

\subsection{Forwards avec dividendes discrets}
Le temps est mesuré en années à partir de la date de valorisation. Soit $S$ le spot, $r(t)$ le taux zéro-coupon en
composition continue utilisé pour l'actualisation et le financement, et $b(t)$ la courbe zéro-coupon de \emph{repo}
(coût d'emprunt du titre), de sorte que détenir l'action financée par un repo rapporte $r-b$. On note $R(t)=r(t)\,t$,
$B(t)=b(t)\,t$ et $G(t)=e^{R(t)-B(t)}$. Des dividendes en numéraire $d_i$ sont détachés aux dates $t_i$.

\begin{definition}[Forward]\label{def-forward}
Le forward de livraison $T$ vaut
\begin{equation}\label{eq-forward}
F(T)=G(T)\bigl(S-D(T)\bigr),\qquad D(T)=\sum_{t_i\le T}\frac{d_i}{G(t_i)}.
\end{equation}
$D(T)$ est le \emph{dividende cumulé actualisé au portage} : les dividendes sont actualisés au taux de portage $r-b$.
\end{definition}

\Cref{eq-forward} est l'argument de réplication habituel : acheter l'action financée à $r-b$ et reporter les
dividendes. Avec des courbes plates, $F(T)=Se^{(r-b)T}-\sum_{t_i\le T}d_ie^{(r-b)(T-t_i)}$. On note que $D$ dépend de la
courbe de repo à travers $G(t_i)$ ; ce couplage est faible (quelques points de base de repo par an sur l'actualisation
des dividendes), mais le bootstrap ci-dessous le traite exactement.

\subsection{L'espace de portage}
\begin{definition}[Portage]\label{def-carry}
Le \emph{repo cumulé} ou \emph{portage} est $B(T)=b(T)\,T$. Pour une grandeur observée aux piliers $T_1<\dots<T_K$, on
note $B_k=B(T_k)$, $b_k=b(T_k)$, $D_k=D(T_k)$, $F_k=F(T_k)$.
\end{definition}

\begin{lemma}[Additivité en portage]\label{lem-additivity}
$\ln F(T)=R(T)-B(T)+\ln\bigl(S-D(T)\bigr)$. Par conséquent, des perturbations du forward, des dividendes et du portage
vérifient, au premier ordre,
\begin{equation}\label{eq-additivity}
\delta B(T)=-\,\delta\ln F(T)-\frac{\delta D(T)}{S-D(T)}+\delta R(T).
\end{equation}
\end{lemma}
\begin{proof}
On prend le logarithme de \cref{eq-forward} et on différentie, avec $\delta\ln(S-D)=-\delta D/(S-D)+O(\delta D^2)$.
\end{proof}

\begin{insight}
Dans l'espace de portage, les mécanismes s'additionnent et rien n'est amplifié. Une erreur de portage $\varepsilon$
à la maturité $T$ est une erreur de repo $\varepsilon/T$ : la même erreur de $0{,}1\,\%$ sur le forward vaut $120$\,pb de
repo à un mois et $2$\,pb à cinq ans. C'est pourquoi les courbes de repo implicite paraissent si bruitées sur le court
terme, et pourquoi tous les tests de cette note portent sur $B$ et non sur $b$.
\end{insight}

\subsection{L'application repo implicite}
Étant donnés des forwards $F_1,\dots,F_K$ aux piliers, la courbe des taux et les dividendes, le \emph{repo implicite}
est la courbe de repo, affine par morceaux en $B$ (repo forward constant entre piliers, $B(0)=0$), qui reprice les
forwards. On la résout pilier par pilier. On fixe $B_1,\dots,B_{k-1}$ et on pose $x=B_k$. Pour un dividende situé entre
le pilier précédent et $T_k$, $B(t_i)=B_{k-1}+w_i(x-B_{k-1})$ avec $w_i=(t_i-T_{k-1})/(T_k-T_{k-1})\in(0,1]$, de sorte que
le forward du modèle en $T_k$ vaut
\begin{equation}\label{eq-phik}
\Phi_k(x)=Se^{R_k-x}-\sum_{t_i\le T_{k-1}}d_i\,e^{(R_k-R_i)-(x-B(t_i))}
-\sum_{T_{k-1}<t_i\le T_k}d_i\,e^{(R_k-R_i)-(1-w_i)(x-B_{k-1})}.
\end{equation}

\begin{lemma}[Monotonie]\label{lem-monotone}
$\Phi_k'(x)\le-\Phi_k(x)$ pour tout $x$. En particulier $\Phi_k$ est strictement décroissante là où le forward est
positif, et $g_k(x)=\ln\Phi_k(x)-\ln F_k$ vérifie $g_k'(x)\le-1$.
\end{lemma}
\begin{proof}
En dérivant \cref{eq-phik} terme à terme,
$\Phi_k'(x)=-Se^{R_k-x}+\sum_{t_i\le T_{k-1}}d_ie^{(\cdot)}+\sum_{T_{k-1}<t_i\le T_k}(1-w_i)\,d_ie^{(\cdot)}$.
Chaque terme de dividende est positif et $1-w_i\le1$ ; remplacer $(1-w_i)$ par $1$ ne peut qu'augmenter le membre de
droite, qui vaut alors $-\Phi_k(x)$. En divisant par $\Phi_k>0$, on obtient $g_k'\le-1$.
\end{proof}

\begin{proposition}[Bootstrap bien posé]\label{prop-bootstrap}
On suppose $S-D(T_{k-1})>0$ pour les dividendes déjà résolus. Alors pour tout $F_k>0$ il existe un unique $B_k$ tel que
$\Phi_k(B_k)=F_k$ et $\Phi_k>0$ sur la branche de la solution. La méthode de Newton appliquée à $g_k$ converge
localement de façon quadratique, et globalement si on la sécurise par dichotomie.
\end{proposition}
\begin{proof}
Quand $x\to-\infty$,
\[
e^{x}\Phi_k(x)\to e^{R_k}\Bigl(S-\sum_{t_i\le T_{k-1}}d_ie^{-R_i+B(t_i)}\Bigr)=e^{R_k}\bigl(S-D(T_{k-1})\bigr)>0,
\]
car les termes de la période courante ne croissent que comme $e^{-(1-w_i)x}$ ; donc $\Phi_k\to+\infty$. Quand $x$ croît,
$\Phi_k$ décroît strictement tant qu'elle est positive (\cref{lem-monotone}) et soit tend vers $0$, soit change de
signe ; par continuité, elle prend chaque valeur de $(0,\infty)$ exactement une fois sur la branche où elle est
positive. Sur cette branche, $g_k$ est $C^2$ avec $|g_k'|\ge1$, donc l'itération de Newton $x\mapsto x-g_k(x)/g_k'(x)$
converge quadratiquement au voisinage de la racine ; la sécurité par dichotomie (conserver un encadrement $[\ell,u]$
avec $g_k(\ell)>0>g_k(u)$ et revenir à son milieu si nécessaire) rend la convergence globale.
\end{proof}

\begin{example}[Un bootstrap à la main]\label{ex-bootstrap}
Spot $100$, taux plat de $3\,\%$, dividendes de $1{,}5$ en $t=0{,}30$ et de $2{,}0$ en $t=0{,}80$, piliers $0{,}25$,
$0{,}5$ et $1$, repo vrai de $80$, $70$ et $60$\,pb. Les forwards et le repo implicite sont
\begin{center}\input{generated/\lang/tab_boot}\end{center}
Au troisième pilier, le résidu de Newton $|g_3|$ passe de \vNewtonA{} à \vNewtonB{} en une itération, puis à zéro à
la suivante. Sans dividende, $g_k$ est affine et une seule itération suffit.
\end{example}

\subsection{Forwards à partir des prix d'options de consensus}
\begin{lemma}[Parité et forward de consensus]\label{lem-parity}
Supposons que chaque contributeur $j$ cote des calls et des puts européens avec un facteur d'actualisation commun $DF$.
Alors
\begin{enumerate}[label=(\roman*)]
\item $C_j(K)-P_j(K)=DF\,(F_j-K)$ pour tout strike ;
\item si le consensus moyenne les mêmes contributeurs à chaque strike,
$\overline{C}(K)-\overline{P}(K)=DF\,(\overline{F}-K)$, où $\overline{F}$ est la moyenne des forwards des contributeurs ;
\item le strike $K^\ast$ où $\overline{C}=\overline{P}$ est égal à $\overline{F}$ et ne dépend pas de $DF$ ;
multiplier tous les prix par un même facteur ne déplace pas $K^\ast$.
\end{enumerate}
\end{lemma}
\begin{proof}
(i) est la parité call--put~\cite{hull}. (ii) s'obtient en moyennant (i), qui est linéaire. (iii) On annule (ii) ;
multiplier les deux prix par un facteur multiplie la droite $K\mapsto\overline{C}-\overline{P}$ sans déplacer sa racine.
\end{proof}

\begin{example}\label{ex-parity}
Avec des strikes $90$, $100$ et $110$, des puts à $2{,}1$, $5{,}4$ et $10{,}9$ et des calls construits à partir de
$F=102{,}3$ et $DF=0{,}97$, la régression pondérée de $\overline{C}-\overline{P}$ sur $K$ donne $F=\vParF$ et
$DF=\vParDF$. Multiplier tous les prix par $1{,}02$ (un facteur d'actualisation erroné) donne $F=\vParFtwo$ et
$DF=\vParDFtwo$ : le forward est inchangé.
\end{example}

\begin{remark}
Deux situations mettent (ii)--(iii) en défaut : un filtrage des valeurs aberrantes appliqué séparément par strike et
par type d'option (les contributeurs moyennés diffèrent alors entre calls et puts), et l'exercice américain (la parité
devient une inégalité). Le $DF$ ajusté, comparé à la courbe, donne un spread de financement implicite qui signale ces
échéances. Le contrôleur conserve sa propre courbe des taux des deux côtés de la comparaison, de sorte que les taux
s'éliminent dans l'écart Totem--FO.
\end{remark}

\subsection{Périodes de dividendes et allocation}
\begin{definition}[Allocation]\label{def-allocation}
Les niveaux de dividendes Totem $L_y$ sont des sommes de dividendes détachés dans la période $y=(e_{y-1},e_y]$
(en général des années civiles). Un profil attribue des poids $\pi_i\ge0$ avec $\sum_{i\in y}\pi_i=1$ aux dates de
détachement de chaque période ; les dividendes alloués sont $d_i=L_y\pi_i$. Le contrôleur prend normalement le
calendrier de détachement du FO comme profil.
\end{definition}

\begin{lemma}\label{lem-allocation}
À $G$ fixé, $D_k=\sum_yA_{ky}L_y$ avec $A_{ky}=\sum_{i\in y,\,t_i\le T_k}\pi_i/G(t_i)$.
\end{lemma}
\begin{proof}
On substitue $d_i=L_y\pi_i$ dans \cref{eq-forward} et on regroupe par période.
\end{proof}
$A_{ky}$ est la part, actualisée au portage, de la période $y$ détachée avant l'échéance $k$. Elle vaut $0$ ou
environ $1/G$, sauf pour les périodes qui chevauchent une échéance : c'est exactement là que le profil compte.

\subsection{Modèle des données Totem}
\begin{assumption}[Contributeurs cohérents]\label{ass-consistent}
Le contributeur $j$ a des vues $b_j(\cdot)$ et $D_j(\cdot)$ et cote des prix d'options cohérents avec
$F_j(T)=G_j(T)(S-D_j(T))$, avec la courbe des taux et le spot communs.
\end{assumption}
\begin{assumption}[Panels]\label{ass-panels}
À l'échéance $T_k$, le panel options $\Aset_k$ compte $n_{A,k}$ contributeurs ; pour la période $y$, le panel
dividendes $\Bset_y$ en compte $n_{B,y}$. Les panels sont emboîtés en maturité ($\Aset_{k+1}\subseteq\Aset_k$,
$\Bset_{y+1}\subseteq\Bset_y$). Les erreurs de vue sont indépendantes d'un contributeur à l'autre. Les deux services ont
$m$ contributeurs en commun ($m=0$ pour des panels disjoints).
\end{assumption}
\begin{assumption}[Consensus]\label{ass-consensus}
Le consensus est la moyenne arithmétique sur le panel. Totem publie, par échéance et par période, le consensus, la
dispersion entre contributeurs et leur nombre.
\end{assumption}

\begin{definition}[Repo implicite Totem et écart de portage]\label{def-implied}
Le repo implicite Totem $b^\T$ est le bootstrap des forwards de consensus $\overline{F}_k$ avec les dividendes de
consensus alloués $\overline{L}_y$. Avec le repo FO $b^\FO$ aux mêmes piliers, l'\emph{écart de portage} est
\begin{equation}\label{eq-y}
y_k=B^\T_k-B^\FO_k=\bigl(b^\T_k-b^\FO_k\bigr)\,T_k,\qquad k=1,\dots,K.
\end{equation}
\end{definition}
